\section{Method}\label{sec:method}

\subsection{Problem Setting}

We are given a distribution $p(\mathcal{T})$ over tasks.
Each task $\tau \sim p(\mathcal{T})$ defines a joint data distribution $p_\tau(x, y)$.
A \emph{support set} $S = \{(x_i, y_i)\}_{i=1}^N$ of $N$ i.i.d.\ draws is provided at meta-test time.
The goal is to produce the parameters $w$ of a small \emph{target network} $h_w : \mathcal{X} \rightarrow \mathcal{Y}$ that generalises from $S$ to held-out query points $(x, y) \sim p_\tau$.
The goal is deterministic adaptation from task evidence to effective task-specific weights.

\subsection{Architecture}

The main system used in our study has two core components: an encoder $f$ and a hypernetwork decoder $g$.

\paragraph{Attention set encoder.}
The encoder maps a variable-length support set $S$ to a pair $(c, z)$ where $c \in \mathbb{R}^{d_c}$ is a \emph{context} vector and $z \in \mathbb{R}^{d_z}$ is a \emph{latent code}.
Each support pair $(x_i, y_i)$ is projected by a 2-layer MLP to a 128-dimensional token.
Three Transformer-encoder layers (4 attention heads, hidden dimension 128, GELU activations, LayerNorm) process the token sequence.
An attention-pooling layer with a learned query vector aggregates the sequence into a single vector, which is then projected to $c \in \mathbb{R}^{64}$ and $z \in \mathbb{R}^{32}$ by separate MLPs.

\paragraph{Hypernetwork decoder.}
A three-layer MLP with hidden dimension 256 and SiLU activations maps the concatenation $[z; c] \in \mathbb{R}^{96}$ to a flat parameter vector, which is reshaped to the weight tensors of the target network.
The target network is a 4-layer ReLU MLP with 64 hidden units (input dimension 2, output dimension 1, approximately 13K parameters).

\paragraph{Text encoder.}
For text-conditioned encoding, task descriptions are processed by DistilBERT (pre-trained, frozen) to produce a 768-dimensional CLS token embedding.
A two-layer MLP projects this to $[c; z] \in \mathbb{R}^{96}$, matching the support encoder output space.
For hybrid mode, the context-latent pairs from support and text branches are mixed as $\alpha \cdot [c_{\text{sup}}; z_{\text{sup}}] + (1-\alpha) \cdot [c_{\text{txt}}; z_{\text{txt}}]$ with $\alpha = 0.5$.
\paragraph{Oracle: task identification baseline.}
For the oracle modality, we provide a one-hot encoded family index directly as input.
If task $\tau$ belongs to family $j \in \{1, \ldots, F\}$, the oracle input is a one-hot vector $\mathbf{e}_j$.
A two-layer MLP projects this $F$-dimensional input to $[c; z] \in \mathbb{R}^{96}$.
Oracle represents perfect task identification with zero structural information: the model knows exactly which family the task belongs to, but knows nothing about the task's input-output relationships, dynamics, or semantic intent.
This modality serves as a diagnostic test: if oracle performs well, it would suggest task labels alone suffice for decoder generalization; if oracle fails, it would indicate that structural task information is necessary.
\subsection{Training}

Training proceeds in two sequential stages.

\paragraph{Stage 1: encoder and hypernetwork.}
We jointly optimize the encoder, hypernetwork decoder, text projector, and target network.
For each episode, a batch of $B{=}32$ tasks is sampled; each task provides $N{=}16$ support points.
The encoder (or text/oracle branch) produces $(c, z)$; the decoder maps $(z, c)$ to target weights $w$; the target network evaluates on held-out query points.
The loss is task-prediction MSE (regression/control) or cross-entropy (classification).
This stage trains the decoder to extract task structure from the encoded representation and predict unseen query points---capturing how the decoder converts task information into task-specific weights.
Optimiser: AdamW with learning rate $10^{-3}$, 1000 training steps.

\subsection{Inference and Evaluation}

At meta-test time, given support set $S$, we:
\begin{enumerate}
  \item Run the encoder (or text/hybrid branch) to obtain context $c$.
  \item Use the deterministic encoder latent $z$ to decode target weights.
  \item Evaluate the resulting target networks on held-out query data.
\end{enumerate}
The FW (det.) path is the primary model in the paper.

\paragraph{No-adapt MLP vs adaptive learners.}
The \emph{no-adapt MLP} is still a learned model, but it is trained as a single
task-agnostic predictor and does not condition on the support set at test time.
By contrast, FW (det.) and Direct (no FW) systems are \emph{adaptive}: they ingest
support data for each episode and produce task-conditioned predictions.
Thus, reward comparisons against the no-adapt MLP measure the marginal value of
per-episode adaptation over a strong non-adaptive learned policy.
